\subsection{Explaining a recommendation}
\label{sec:res-explanation}

The record of Step~5 allows any recommendation to be traced back to the operation. Fig.~\ref{fig:explanation} shows this for the takeover requests of S1 and S2, which reached the threshold through different parts of the network; the numbers in parentheses refer to the rows of Table~\ref{tab:arc-trace}. In S1, the request followed the collision pathway. \node{SonarRange} Critical raised the probability of a likely \node{CollisionWithObstacle} from 0.02 to 0.92. The arc comes from rule A5: a return in the safety box is a pivotal event on the path to collision in the transit ESD (26). \node{CollisionWithObstacle} is in turn the collision pathway to the takeover constraint C1 (A5, N7; 39), and the probability of \emph{TakeoverRequested} rose from 0.14 to 0.68, while the other parents of the output did not change. In S2, \node{CollisionWithObstacle} stayed at 0.12, and the request followed the navigation pathway. \node{USBLStrength} and \node{WaypointTrackingError} are read for the navigation-quality decision of the inspection CoTA (AI1.1 $\rightarrow$ AI1.2; A1, A2; 15, 18), and \node{Current} degrades position keeping (N4; 17). Together they raised the probability of an inadequate \node{NavigationSystemPerformance} from 0.01 to 0.60 at the request, and to 0.92 one update later. A valid position is required for adequate inspection data (I-2, I-4 $\rightarrow$ I-6 $\rightarrow$ I-X3; A5; 34), so the probability of an inadequate \node{DataQuality} rose to 0.49 and then 0.73; this pathway reaches the output only in inspection (A5, N7, A4; 40, 41). \node{AutonomousControlPerformance} rose with it, to 0.58 and then 0.85, and reaches the output through the request-and-confirm interaction A-TO.3 $\rightarrow$ H-TO.1 (A1, A3; 30, 37). In neither scenario did the supervisor branch contribute: the probability of a low \node{HumanDecisionReliability} stayed between 0.13 and 0.15. The supervisor therefore sees not only that a takeover is requested but which task of the vehicle has failed: obstacle clearance in transit (AT3) in S1, and the navigation-quality decision on which the inspection data depend (AI1.2) in S2. The two call for different responses, steering clear of the obstacle in one case and re-establishing the position before the inspection continues in the other, which is the information needed to take over in an informed way.

\begin{figure*}
    \centering
    \includegraphics[width=\textwidth]{figures/results/explanation_example.pdf}
    \caption{Explanation of the takeover requests in (a)~S1 and (b)~S2 (the open-loop runs of Fig.~\ref{fig:sim-traces}). Evidence nodes (left) are coloured by their state at the request; the other nodes give the probability of their worst state before the injection $\rightarrow$ at the request, and the output node also gives $U$. Black: the path that carried the request; grey: nodes and arcs that did not change. Each arc on the path is labelled with its mapping rule (Table~\ref{tab:mapping-rules}) and its row in Table~\ref{tab:arc-trace}, whose source in the task analysis is listed under each panel. Abbreviations as in Fig.~\ref{fig:gating}.}
    \label{fig:explanation}
\end{figure*}
